\section{Admissible unit actions and a cyclotomic trace invariant}
\label{sec:admissible-arrows}

This section uses the actual all-open-subgroup condition in
\cite[Example~1.8(iv), author-version pp.~38--39]{IUT2}.
It must be distinguished from allowing every automorphism of one fixed
integral lattice. Fix a finite extension $E/\Q_p$, an algebraic closure
$\overline E$, and $G=\operatorname{Gal}(\overline E/E)$. For each finite
$K/E$ inside $\overline E$, put
\[
 U_K=\OO_K^\times,\qquad L_K=\log U_K\subset K,
 \qquad I_K=p^{-1}L_K.
\]
The logarithm is used only on units. Its kernel on $U_K$ is the finite
group of roots of unity in $K$, and $L_K$ is a full $\Z_p$-lattice in $K$.
These assertions hold also for $p=2$.

Logarithm induces a $G$-equivariant additive isomorphism
\begin{equation}\label{eq:unit-log-quotient}
 \OO_{\overline E}^\times/\mu(\overline E)
       \simeq (\overline E,+).
\end{equation}
Indeed the finite-field kernel statement proves injectivity. Given
$x\in\overline E$, choose $n$ large enough for $\exp(p^n x)$ to converge
in a finite extension and choose a $p^n$-th root $y$ of it. Then $y$ is
a unit and $\log y=x$, proving surjectivity. If
$H=\operatorname{Gal}(\overline E/K)$, the invariants of the quotient
in \eqref{eq:unit-log-quotient} are $K$, whereas the image of the
invariants of the original unit group is precisely $L_K$.

Under this identification, the group $\operatorname{Ism}(G)$ specified
in \cite{IUT2} consists of additive $G$-equivariant automorphisms $F$ of
$\overline E$ preserving $L_K$ for \emph{every} finite $K/E$. The source
requires preservation of these lattices, not only of their rational
spans. Such a map is $\Q_p$-linear on each $K$: preservation of
$p^nL_K$ implies continuity, which extends integer linearity to
$\Z_p$-linearity and then to $\Q_p$-linearity.

\subsection{The all-open condition forces scalar action}

\begin{lemma}[Realizing logarithmic sublattices by norms]
\label{lem:norm-sublattice}
For every finite-index $\Z_p$-submodule $J\subset L_E$, there is a
finite abelian extension $K/E$ such that
\[
 \operatorname{Tr}_{K/E}(L_K)=J.
\]
\end{lemma}
\begin{proof}
Let $V=\log^{-1}(J)\subset U_E$, choose a uniformizer $\pi$ of $E$, and
put $N=\pi^{\Z}V$. This is an open finite-index subgroup of $E^\times$.
The local existence theorem of class field theory
\cite[Chapter~I, Theorem~1.4]{MilneCFT} supplies a finite abelian $K/E$
with $\operatorname{Norm}_{K/E}(K^\times)=N$. For integer-normalized
valuations,
\[
 v_E(\operatorname{Norm}_{K/E}z)=f_{K/E}v_K(z).
\]
Consequently a norm is a unit exactly when its preimage is a unit, and
$\operatorname{Norm}_{K/E}(U_K)=N\cap U_E=V$.
The identity $\log\operatorname{Norm}z=\operatorname{Tr}\log z$ follows
by summing the logarithms of the conjugates. Therefore
$\operatorname{Tr}_{K/E}(L_K)=\log V=J$.
\end{proof}

\begin{lemma}[Sublattice rigidity]\label{lem:sublattice-rigidity}
An automorphism of a nonzero finite free $\Z_p$-module preserving every
finite-index submodule is multiplication by one element of $\Z_p^\times$.
\end{lemma}
\begin{proof}
Choose a basis $e_1,\ldots,e_d$. Preservation of
$\Z_pe_i+p^nL$ for every $n$ puts the image of $e_i$ in their
intersection, which is $\Z_pe_i$: every other coordinate is divisible
by every power of $p$ and hence is zero. Write $F(e_i)=\lambda_i e_i$.
The same argument for $e_i+e_j$ implies $\lambda_i=\lambda_j$.
Thus $F=\lambda\operatorname{id}$; surjectivity shows that $\lambda$
is a unit. The rank-one case is already covered by the first step.
\end{proof}

\begin{theorem}[Classification of the specified Ism action]
\label{thm:ism-scalar}
For the group defined in \cite[Example~1.8(iv)]{IUT2},
\[
 \operatorname{Ism}(G)=\Z_p^\times
\]
under \eqref{eq:unit-log-quotient}, with the right side acting by scalars.
\end{theorem}
\begin{proof}
For finite abelian $K/E$, equivariance gives
$F_E\operatorname{Tr}_{K/E}=\operatorname{Tr}_{K/E}F_K$.
Since $F_K(L_K)=L_K$, Lemma~\ref{lem:norm-sublattice} implies that
$F_E$ preserves every finite-index submodule of $L_E$.
Lemma~\ref{lem:sublattice-rigidity} gives
$F_E=\lambda_E\operatorname{id}$ with $\lambda_E\in\Z_p^\times$.
Repeat over any finite $K/E$: the hypothesis includes all finite
extensions of $K$, since their Galois groups are open in $G$.
Thus $F_K=\lambda_K\operatorname{id}$. Restriction to the nonzero
subfield $E$ gives $\lambda_K=\lambda_E$, proving the assertion on
all of $\overline E$. Conversely every $\Z_p$-unit scalar preserves
all the lattices and commutes with $G$. It is the image of a unit of
$\widehat\Z$ under the unit-power action; prime-to-$p$ components act
trivially after quotienting by all roots of unity.
\end{proof}

Only preserving $L_E$ would leave the full group
$\GL_{[E:\Q_p]}(\Z_p)$; the finite-extension hypothesis is decisive.
The full-unit lifting theorem \cite[Theorem~7.6]{Hoshi} by itself does
not classify automorphisms of the torsion quotient. The preceding proof
supplies that additional step. For the indeterminacy called Ind2 in
\cite[Theorem~3.11(i)]{IUT3}, the theorem applies to its indicated Ism
factors. It does not classify the Galois arrows of Ind1 or the set
operations of Ind3.

For example, if $1<e(E/\Q_p)\le p-2$, logarithm and exponential on the
maximal ideal give $I_E=\pi^{1-e}\OO_E$. The unit
$u=\log(1+p)/p$ has native valuation zero, and all its Ism images
still have valuation zero. None reaches the least valuation $1/e-1$
of this larger lattice. This statement concerns Ism; it adds no
unproved restriction to Ind1.

\subsection{What the procession retains, and an exact Ind2 hull equality}

The Ind1 source in \cite[Theorem~3.11(i)]{IUT3} is the mono-analytic
procession of \cite[Proposition~6.9(ii)]{IUT1}. At a nonarchimedean
place its retained category is $\mathcal B(E)^0$, the connected finite
etale covers of $E$, rather than the earlier curve-covering category
\cite[Examples~3.2(i), 3.3(i), Definition~4.1(iii)--(iv)]{IUT1}.
Its self-equivalences, up to natural isomorphism, correspond to
$\operatorname{Out}(G_E)$: pullback of finite transitive $G_E$-sets
constructs one direction, and adjoining finite disjoint unions and
recovering the fibre functor constructs the inverse.

For fixed capsule labels, the capsule-full poly-isomorphism contains
all collections of constituent isomorphisms
\cite[\S0, pp.~33--34; Definition~4.10]{IUT1}. Consequently every
local outer Galois automorphism occurs as a representative in an
allowed poly-automorphism, with the same representative used at repeated
labels if required. To see this directly, identify each labelled
category with $\mathcal B(E)^0$, conjugate the chosen pullback
equivalence to each copy, and use identity maps at the other places.
Keep the capsule and stage injections equal to the identity. Each
resulting collection belongs to the required full set of constituent
isomorphisms, and conjugation preserves the full connecting sets.
There is no further requirement to lift back to the curve category.
This is a statement about representatives inside poly-morphisms, not
a claim that each representative is a distinct object or morphism of
the coarse category. It supplies no arbitrary integral linear lift.

\begin{proposition}[Ind2 does not enlarge a fixed module hull]
\label{prop:ind2-hull-equality}
On a fixed marked nonarchimedean tensor carrier, let
\[
 T=\bigotimes_{i,\Q_p}\Bigl(\prod_{v\in J_i}E_{i,v}\Bigr),
 \qquad
 A=\bigotimes_{i,\Z_p}\Bigl(\prod_{v\in J_i}\OO_{E_{i,v}}\Bigr)
 \subset B,
\]
where $B$ is the maximal order of the finite etale algebra $T$.
For any subset $S\subset T$, write $\Gamma_2S$ for the union of its
images under the local Ind2 operators of \cite[Theorem~3.11(i)]{IUT3}.
Then
\[
 \Span_B(\Gamma_2S)=\Span_B(S).
\]
The same equality holds for closed $B$-module hulls.
\end{proposition}
\begin{proof}
Theorem~\ref{thm:ism-scalar} makes the action on each local direct
summand multiplication by a scalar $\lambda_{i,v}\in\Z_p^\times$.
In factor $i$, their tuple $c_i=(\lambda_{i,v})_v$ is a unit of
$\prod_v\OO_{E_{i,v}}$. The tensor product operator is multiplication
by $c=\bigotimes_i c_i\in A^\times$, with inverse the tensor of the
$c_i^{-1}$. Thus every image of $S$ lies in $\Span_B(S)$.
The identity Ind2 operator supplies the reverse inclusion of spans.
Taking closures gives the closed version. Any extra convention forcing
some scalars to coincide leaves the argument unchanged.
\end{proof}

This removes the Ind2 step from a $B$-module-hull computation once
the Ind1 image set and native markings are fixed. It does not remove
the distinction between $A$ and $B$, turn an Ind1 map into a
$B$-linear map, or identify Ind3 with multiplication by a unit.

\subsection{Integral Galois transport preserves trace depth}

A coefficient-compatible Galois arrow consists of a continuous
isomorphism $\alpha:G_{E'}\simeq G_E$ and an isomorphism
$\beta:\alpha^*\Z_p(1)_E\simeq\Z_p(1)_{E'}$ of integral Tate modules.
The cyclotomic character and inertia are group-theoretically determined
\cite[Propositions~1.1--1.2 and Corollary~1.3]{MochLocal}.
In particular the coefficient isomorphism can be chosen integrally;
one cannot insert multiplication by $p$ as an isomorphism.

\begin{proposition}\label{prop:trace-transport}
The induced logarithmic map $F_{\log}:E\longrightarrow E'$ satisfies
\begin{equation}\label{eq:trace-transport}
 \operatorname{Tr}_{E'/\Q_p}(F_{\log}x)
       =c_\alpha\operatorname{Tr}_{E/\Q_p}(x)
 \quad(x\in E)
\end{equation}
for some $c_\alpha\in\Z_p^\times$. The statement also holds with both
logarithmic coordinates divided by $p$.
\end{proposition}
\begin{proof}
Put $h_E=\log_p\chi_E\in H^1(G_E,\Z_p)$, where the latter coefficient
action is trivial. Equivariance of $\beta$ gives
$\chi_E\alpha=\chi_{E'}$, hence $\alpha^*h_E=h_{E'}$.
Local duality and the invariant map identify
$H^2(G_E,\Z_p(1))$ with $\Z_p$. Integral cohomology transport is an
isomorphism, so it acts there by a unit $c_\alpha$, not by an arbitrary
nonzero rational $p$-adic scalar.

The local reciprocity cup-product formula, for an integral character $h$
and a Kummer class $\kappa(z)$, is
\[
 \operatorname{inv}(h\mathbin\cup\kappa(z))=h(\operatorname{rec}_E z).
\]
It follows at finite level from the cyclic-algebra description and then
by inverse limit; see \cite[Chapter~III, Proposition~3.6 and \S4]{MilneCFT}.
In the arithmetic-Frobenius convention,
\[
 h_E(\operatorname{rec}_E z)
   =-\log_p\operatorname{Norm}_{E/\Q_p}z
   =-\operatorname{Tr}_{E/\Q_p}\log_p z
 \qquad(z\in U_E).
\]
The integral unit Kummer submodule is also preserved: it is the
annihilator, under this pairing, of the unramified $\Z_p$-character.
Characteristic inertia carries a generator of that unramified
character module to a unit multiple of a generator, and cup-product
naturality preserves its annihilator.

Applying the same naturality with $h_E$ now proves
\eqref{eq:trace-transport} on $L_E$. Torsion unit classes are killed
by logarithm and pair to zero in the torsion-free integral $H^2$;
the argument therefore descends to the actual free unit lattice.
Since $L_E$ spans $E$, $\Q_p$-linearity proves the formula everywhere.
Uniform division by $p$ on both sides does not change it.
\end{proof}

Composing with the Ism factors of Theorem~\ref{thm:ism-scalar} only
changes $c_\alpha$ by a unit. Thus the valuation of a nonzero absolute
trace is an invariant of these local Galois/Kummer/Ism arrows.
For example, on $E=\Q_7(\pi)$ with $\pi^3=7$, the normalized integral
log lattice $I_E$ has basis $1,\pi/7,\pi^2/7$. Swapping its first two basis vectors is
not an admissible Galois/Kummer arrow: their traces are $3$ and $0$.
No claim that every matrix satisfying \eqref{eq:trace-transport} is
admissible is made.

\begin{theorem}[A native root-power obstruction]\label{thm:root-trace}
Let $p\ne\ell$ be odd primes, $b\in\Q_p$ with
$N=v_p(b)>0$ and $\ell\nmid N$, and let $E/\Q_p$ contain
$\mu_\ell$ and $a$ with $a^\ell=b$. Write $d=[E:\Q_p]$ and, for
$s\ge1$ with $\ell\nmid s$, put $\tau_s=\log(1+pa^s)/p$. Then
\begin{align}
 \operatorname{Tr}_{E/\Q_p}\tau_s
   &=\frac{d}{\ell p}\log(1+p^\ell b^s),\label{eq:root-trace}\\
 v_p(\operatorname{Tr}_{E/\Q_p}\tau_s)
   &=v_p(d/\ell)+\ell-1+sN.\label{eq:root-trace-depth}
\end{align}
In particular, distinct such $s$ give distinct orbits under the arrows
of Proposition~\ref{prop:trace-transport}, even after Ism action.
\end{theorem}
\begin{proof}
The field $F=\Q_p(\mu_\ell)$ is unramified over $\Q_p$.
The valuation assumption makes $b$ a non-$\ell$-th-power in $F$, so
$K=F(a)$ has degree $\ell$ over $F$. Since multiplication by $s$
permutes the exponents modulo $\ell$, and $\ell$ is odd,
\[
 \operatorname{Norm}_{K/F}(1+pa^s)
   =\prod_{\zeta\in\mu_\ell}(1+p\zeta a^s)=1+p^\ell b^s.
\]
All factors are principal units. Taking logarithms, dividing by $p$,
and using trace transitivity gives \eqref{eq:root-trace}: the remaining
factor is $[E:K][F:\Q_p]=d/\ell$.
For $x=p^\ell b^s$, the term $x$ is the unique least-valuation term
in $\log(1+x)$, since
\[
 v_p(x^k/k)-v_p(x)=(k-1)(\ell+sN)-v_p(k)>0\quad(k\ge2).
\]
This proves \eqref{eq:root-trace-depth}, including nonvanishing.
Its values differ by $(s-t)N$ for distinct labels, which
\eqref{eq:trace-transport} cannot change.
\end{proof}

For a concrete example take $p=11$, $\ell=5$, $b=11$, and
$E=\Q_{11}(\pi)$ with $\pi^5=11$. The base field contains $\mu_5$;
the trace depths of $\tau_1$ and $\tau_4$ are $5$ and $8$.
This is a specified local example, not a construction of all the
global initial theta data.

There is a direct local link with a split Tate curve having rational
full $2$-torsion. If $q$ is its parameter, choose $b^2=q$.
The class of $b$ in $\overline{\Q}_p^\times/q^\Z$ is rational, so
$g(b)/b\in q^\Z\cap\{1,-1\}$. Valuation makes this intersection
$\{1\}$; hence $b\in\Q_p$. Rationalizing full $\ell$-torsion also
rationalizes full $2\ell$-torsion and supplies the required $a$ and
$\mu_\ell$. The theorem consequently applies when
$\ell\nmid v_p(q)/2$. It does not assert that every bad prime meets
that condition.

For a pure tensor with $m-1$ factors $u=\log(1+p)/p$ and one factor
$\tau_s$, trace in $E^{\otimes m}$ is the product of the factor traces.
Its valuation is
\[
 (m-1)v_p(d)+v_p(d/\ell)+\ell-1+sN.
\]
Local arrows, Ism actions, and factor permutations preserve this number,
including when labels require repeated use of one arrow.
This separates the exponent-$1$ input from the exponent-$j^2$ input
for $j>1$, $j<\ell$. Raising $1+pa$ to the power $j^2$, which scales
its logarithm, is a different operation from replacing $a$ by $a^{j^2}$.

\subsection{Why this is not yet a holomorphic-hull comparison}

Multiplying generators by coefficients of the maximal tensor order
need not preserve their trace depths. Thus distinct pure orbits do not
imply distinct module hulls. The upper-semicompatibility operation Ind3
in \cite[Theorem~3.11(ii)]{IUT3} is a further operation, and has not
been replaced here by an automorphism of one fixed Kummer module.
Nor does changing an untilt's real-valued norm change the native trace
in a fixed marked $p$-adic field.

A converse test is useful: trace preservation alone permits a common
least-layer displacement, without establishing an Ind1 lift.
\begin{lemma}[Common avoidance with fixed trace]
\label{lem:trace-transvection}
Let $k$ be a field of characteristic different from two, let $V$ be a
$k$-vector space, and let $t,l\in V^*$. Suppose $t(w)=0$, $l(w)=1$.
For any nonzero $x,y\in V$, there is an automorphism $F$ with
$tF=t$ and $l(Fx),l(Fy)\ne0$.
\end{lemma}
\begin{proof}
Set $P(v)=v-l(v)w$. Choose a functional $h$ nonzero on each nonzero
vector among $P(x),P(y)$. Such a functional exists: separate the first
vector; if that functional kills the second, separate the second and
use either the new functional or the sum, according to its value on the
first. Put $f=hP$, so $f(w)=0$. The maps
$F_z(v)=v+zf(v)w$ have inverses $F_{-z}$ and preserve $t$.
For $v=x,y$, the polynomial $l(v)+zf(v)$ is nonzero: if $l(v)=0$,
then $P(v)=v\ne0$. Two nonzero affine polynomials have at most two
roots altogether. The three distinct scalars $0,1,-1$ therefore
contain a valid choice of $z$.
\end{proof}

For finite Galois $E/\Q_p$ with $2\le e\le p-2$, one has
$I_E=\pi^{1-e}\OO_E$, the inverse different. Both
$\operatorname{Tr}:I_E\to\Z_p$ and
$\operatorname{Tr}:\OO_E\to\Z_p$ are onto; the latter follows by
tracing first through the tame ramified extension and then through the
unramified one. On $I_E/pI_E$, trace modulo $p$ is independent of any
nonzero $\mathbb F_p$-coordinate of $I_E/\pi I_E$: the latter vanishes
on $\OO_E\subset\pi I_E$, whereas trace does not. The lemma gives an
automorphism preserving trace modulo $p$ and reaching the least
valuation layer for two primitive vectors simultaneously.
Choose a basis splitting integral trace. The residue matrix has first
row $(1,0,\ldots,0)$; lift that row exactly and the others arbitrarily.
The determinant is a unit, so this is an integral invertible lift
preserving trace exactly. No lift to an absolute Galois automorphism
is established by this construction. Determining that image, and then performing all
source-prescribed hull operations, remains necessary for a global
comparison. No IUT theorem or the abc conjecture has been disproved
by these local nonreachability statements.
